\documentclass[../main.tex]{subfiles}
\begin{document}
\chapter{Exercises Lecture 5}

\section{MC simulation of a 2D Ising model by the Metropolis algorithm}

\noindent\textbf{Program} Write a program that simulates a 2D Ising model with periodic boundary conditions by using the Metropolis acceptance matrix $a_{ij}$ and the matrix $\Gamma_{ij}$ based on the local spin-flip (Glauber) proposed move.

\noindent\textbf{Simulation} By assuming $kB = 1$ and $J = 1$ simulate the 2D Ising model for different values of temperatures (at least 3, one below, one above, and one close to the critical temperature $T_c = \frac{2}{\ln(1+\sqrt{2})}$ and 3 different
values of L (for instance 25,50 and 100).

\noindent\textbf{Equilibration time, averages and fluctuations} After having determined the equilibrium time and disregarding the samples for $t < \tau_{eq}$ estimate the ensemble averages of the magnetization per spin, the energy per spin, and the corresponding fluctuations (specific heat and magnetic susceptibility). 

\noindent \textbf{Integrated correlation time and critical slowing-down} Estimate the autocorrelation time of the magnetization and the energy for the MC simulations proposed above and estimate the errors accordingly.

\subsection{Solution}

The energy of this 2D Ising model given a configuration of the spins $\mathcal{C}$ is:
\begin{align}
H(\mathcal{C}) &= -J \sum_{\langle i, j \rangle} \sigma_i \sigma_j \quad\quad\text{with } \sigma_i=\pm 1\quad \text{and}\quad \langle i, j \rangle\quad\text{it's a sum over n.n.}
\end{align} At every step of the MC, all spins are selected in order and for each we calculate the energy difference if that spin is flipped as:
\begin{align}
\Delta E = 2J\sigma_i \sum_{{\langle j \rangle}_i}\sigma_j  \quad\quad\text{with } {\langle j \rangle}_i\ \text{the 4 neighbors of }\sigma_i
\end{align} If $\Delta E < 0$ or $\exp{(-\Delta E)\beta}$ is bigger than a random number between [0, 1] then we accept the move and flip the spin (this is the Metropolis algorithm).

For each MC steps, we calculate our observables: the energy and magnetization per site. After the system reaches equilibrium, the averages of the observable are calculated, from which we can calculate the specific heat and the magnetic susceptibility.
\begin{align}
    <E_{site}> = \frac{1}{t_{max}-\tau_{eq}}\frac{1}{L^2} \sum_{t=\tau_{eq}}^{t_{max}} H_t
    \quad <m> = \frac{1}{t_{max}-\tau_{eq}}\frac{1}{L^2} \sum_{t=\tau_{eq}}^{t_{max}} M_t
    \quad C_V = \frac{\text{Var[$E_{site}$]}}{k_B T^2}
    \quad \chi_T = \frac{\text{Var[$m$]}}{k_B T}
\end{align}

We also calculated:
\begin{itemize}
    \item $\tau_{eq}$ the time the simulation reached the equilibrium
    \item $\tau_O^{int}$ the integrated autocorrelation time which is:
    \begin{align} 
        \int_{0}^{t_{max}} \frac{C_O(t)}{C_O(0)} dt = \int_{0}^{t_{max}} e^{-t/\tau} dt = \tau_O^{int} = \sum_{t=0}^{t_{max}} \frac{C_O(t)}{C_O(0)}\Delta t
    \end{align}
    which tells us how distant we should sample $O$ to get uncorrelated samples ($C_O(t)$ is the autocorrelation function).
    \item $S_{t_{max}}^2$ which tells us the variance of the mean of the observable $O$ and it is calculated as: \begin{align} 
        S_{O}^2 = \frac{1+2\tau / \Delta t}{(t_{max} - \tau_{eq} - 1)} \sum_{t=\tau_{eq}}^{t_{max}} (O_t - \bar{O}_{t_{max}})^2.
    \end{align}
\end{itemize}

A C++ program calculates the evolution of the spins on the lattice and writes them to a binary file at every step. Then a Python script reads it and analyzes the content.

We run the program with 3 different lattice dimensions (side): 50, 100 and 200, and for 3 different values of temperature: $0.5T_c$, $0.97T_c$ and $1.2T_c$.

The results are shown in table \ref{tab: es_5_1_ising}.

\begin{table}[!ht]
    \centering
    \resizebox{\textwidth}{!}{%
    \begin{tabular}{c| c c c| c c c| c c c}
        \toprule
        \multicolumn{10}{c}{\textbf{Results}}\\
        \midrule
          &\multicolumn{3}{c|}{$T=0.5\ T_c$} & \multicolumn{3}{c|}{$T=0.97\ T_c$} & \multicolumn{3}{c}{$T=1.2\ T_c$}\\
          &$L=50$&$L=100$&$L=200$&$L=50$&$L=100$&$L=200$&$L=50$&$L=100$&$L=200$\\
        \midrule
    $<E_{site}>$    & -1.993 & -1.992   & -1.992   & -1.541& -1.546 & -1.546 & -0.95  & -0.9501  & -0.9501  \\
    $<m>$           & -0.9981& -0.9981  & -0.998   & 0.7735& -0.7845&-0.7863 &$-1.37\cdot 10^{-3}$&$2.15\cdot 10^{-3}$& $2.37\cdot 10^{-4}$\\
    $N_{steps}$     & 2000   & 2000     & 2000     & 16000 & 16000  & 16000  & 8000     & 8000     & 8000     \\
    $C_V$           & 0.04523& 0.04791  & 0.05224  & 1.472 & 1.439  & 1.32   & 0.6411   & 0.5773   & 0.5829   \\
    $\chi_T$        & 0.00363& 0.003803 & 0.004164 & 7.104 & 4.44   & 3.626  & 8.195    & 7.438    & 8.45     \\
    $\tau_{eq}$     & 100    & 100      & 200      & 500   & 4600   & 9600   & 100      & 4200     & 800      \\
    $\tau_E^{int}$  & 1.01   & 1.74     & 0.967    & 30.6  & 29.9   & 14.5   & 29.7     & 5.76     & 8.33    \\
    $\tau_m^{int}$  & 0.967  & 1.03     & 0.857    & 27.6  & 75.6   & 45.9   & 15.6     & 8.04     & 6.95    \\
    $S_{E}^2$       & $7.02 \cdot 10^{-5}$ & $2.76 \cdot 10^{-5}$& $4.93 \cdot 10^{-6}$ & $1.78 \cdot 10^{-1}$ & $4.25 \cdot 10^{-2}$  & $4.8 \cdot 10^{-3}$ & $1.148 \cdot 10^{-1}$& $5.36 \cdot 10^{-3}$ & $1.91 \cdot 10^{-3}$ \\
    $S_{m}^2$         & $4.85 \cdot 10^{-6}$ & $1.32 \cdot 10^{-6}$& $3.21 \cdot 10^{-7}$ & $3.51 \cdot 10^{-1}$ & $1.49 \cdot 10^{-1}$ & $1.851 \cdot 10^{-2}$& $2.87 \cdot 10^{-1}$& $3.46 \cdot 10^{-2}$  & $8.57 \cdot 10^{-3}$ \\
        
        \bottomrule 
     \end{tabular}}
    \caption{Results for Ising 2D}
    \label{tab: es_5_1_ising}
\end{table}

We also show some images for $L=200$ of the $E_{site}$ and $m$ for the three temperatures: figure \ref{fig: ex_5}.

\begin{figure}[!ht]
    \centering
    \begin{subfigure}{.45\textwidth}
        \centering
        \includesvg[width=0.99\columnwidth]{/imgs/ex_5_under_tc.svg}
        % \caption{A subfigure}
        % \label{fig:sub1}
    \end{subfigure}
    \begin{subfigure}{.45\textwidth}
        \centering
        \includesvg[width=0.99\columnwidth]{/imgs/ex_5_near_tc.svg}
        % \caption{A subfigure}
        % \label{fig:sub2}
    \end{subfigure}
    \begin{subfigure}{.45\textwidth}
        \centering
        \includesvg[width=0.99\columnwidth]{/imgs/ex_5_over_tc.svg}
        % \caption{A subfigure}
        % \label{fig:sub2}
    \end{subfigure}
    \caption{Energy and magnetization for three systems with $L=200$ and different temperatures. The red line is when the equilibrium was reached (for both quantities).}
    \label{fig: ex_5}
\end{figure}

On github are available also the animation from $t=0$ to $t\approx \tau_{eq}$: \href{https://github.com/EmanueleSarte/nmsm/tree/main/ex5}{here}.
\end{document}